% ---------------------------------------------------------------------------
% muscle_model_appendix_draft.tex — appendix-ready draft covering the two
% muscle formulations used in this work (MuJoCo muscle actuator as produced
% by the conversion, and OpenSim's Thelen2003/Millard2012 Hill-type models),
% and the quantitative force comparison along the subject01 walking
% trajectory. Written 2026-09-14; regenerate numbers with
% Code\MuJoCo_SNS\spinal\_muscle_force_compare.py (muscle_force_compare.csv,
% figures\muscle_force_compare.png).
% ---------------------------------------------------------------------------

\section{Muscle Model Formulations and Cross-Platform Force Comparison}
\label{sec:muscle-models}

\subsection{MuJoCo Muscle Actuator (as converted)}

The conversion from OpenSim to MuJoCo re-expresses each Thelen2003
muscle as a MuJoCo \texttt{general} actuator with the built-in muscle
dynamics. Activation is a first-order process with distinct rise and
fall time constants (0.01\,s and 0.04\,s in the converted model). Force
is computed from the actuator's normalized length
$\bar{L} = (L - L_{\min}) / (L_{\max} - L_{\min})$, where the length
range is taken from the converted geometry, as
\begin{equation}
  F = F_{\max}\,\bigl[\,a\, f_L(\bar{L})\, f_V(\dot{L}) - f_P(\bar{L})\,\bigr],
  \label{eq:mj-muscle}
\end{equation}
with $a \in [0,1]$ the activation, $f_L$ a force--length curve that
peaks inside the normalized range, $f_V$ a force--velocity curve, and
$f_P$ a passive element. The peak force scale $F_{\max}$ (actuator
parameter \texttt{gainprm[2]}) is set by the converter such that the
peak force over the model's range of motion matches the OpenSim
\texttt{max\_isometric\_force}; audited across all 84 active muscles,
the deviation is $\le 15\%$ (Sec.~\ref{sec:muscle-models-compare}). Two
properties matter for interpretation: (i) the tendon is \emph{rigid} --
there is no series-elastic element and no tendon slack length; force at
the joint equals the contractile-element force immediately; and (ii)
all curve parameters are shared defaults of the converted model rather
than per-muscle fits.

\subsection{OpenSim Thelen2003 and Millard2012 Formulations}

OpenSim's Thelen2003Muscle comprises a contractile element (activation
dynamics as above; active force--length curve; Hill-type force--velocity
curve) in series with a compliant tendon, plus a parallel passive fiber
element, with fiber pennation. Unlike the MuJoCo actuator, the fiber
length is \emph{not} the musculotendon length: at every instant the
model solves the force equilibrium
$F_{\text{tendon}}(L_{\text{mt}} - L_{\text{fiber}}\cos\alpha)
 = F_{\text{fiber}}(L_{\text{fiber}}, \dot{L}_{\text{fiber}}, a)$,
so tendon compliance filters and delays force transmission and stores
energy. The tendon is a quadratic-then-linear ``strip'' curve anchored
at the slack length $L_s$ (force zero until $L_{\text{mt}}$ stretches
the tendon beyond $L_s$). Millard2012EquilibriumModel, OpenSim's newer
default, retains the same equilibrium structure but replaces the curve
shapes (an asymmetric Gaussian active-force--length curve with a
plateau, an exponentially stiffening tendon) and adds fiber damping; it
is numerically better behaved at extreme lengths. Both OpenSim models
therefore differ from the converted MuJoCo plant in the same single
respect: series-elastic compliance.

\subsection{Cross-Platform Force Comparison}
\label{sec:muscle-models-compare}

Both plants were driven with identical kinematics and identical
activations along the subject01 walking trajectory (the same IK frames
used for the static-optimization back-solve). The OpenSim side is the
StaticOptimization solution's own muscle forces (full Thelen dynamics
with compliant tendon); the MuJoCo side replays the identical joint
trajectory with those activations. Table~\ref{tab:force-r2} reports, per
muscle, the coefficient of determination $R^2$ between the two force
time series, the shape correlation, and the best-aligning lag.

\begin{table}[h]
  \centering
  \caption{MuJoCo vs.\ OpenSim muscle force along the subject01 walk at
  matched activations. Lag in analysis frames ($\approx 17$\,ms each).}
  \label{tab:force-r2}
  \small
  \begin{tabular}{lrrr}
    \toprule
    Muscle & $R^2$ & shape corr$^2$ & lag [frames] \\
    \midrule
    soleus   & 0.742 & 0.800 & 1 \\
    med.\ gastrocnemius & 0.888 & 0.894 & 2 \\
    lat.\ gastrocnemius & 0.884 & 0.887 & 2 \\
    tibialis anterior & 0.945 & 0.945 & 0 \\
    vastus lateralis & 0.941 & 0.945 & 1 \\
    rectus femoris & $-1.15$ & 0.384 & 0 \\
    semimembranosus & 0.565 & 0.621 & 0 \\
    biceps femoris short & $-0.26$ & 0.868 & 0 \\
    psoas & 0.927 & 0.934 & 1 \\
    gluteus maximus (2) & 0.926 & 0.943 & 1 \\
    tensor fasciae latae & 0.933 & 0.935 & 0 \\
    peroneus brevis & 0.977 & 0.989 & 0 \\
    \bottomrule
  \end{tabular}
\end{table}

For ten of twelve representative muscles the two formulations agree
with $R^2 = 0.87$--$0.98$ and mean-force ratios within $\pm 22\%$, with
\emph{no} systematic phase shift (best-aligning lag 0--2 frames,
$\le 34$\,ms). The absence of a phase lead in the rigid-tendon plant
indicates that, along physiological trajectories, missing tendon slack
compliance does not desynchronize force transmission; its signature is
instead visible in waveform \emph{shape}: for the soleus (the largest
series-elastic energy store in human gait), the OpenSim compliant
tendon carries substantial force through mid-stance while the MuJoCo
rigid-tendon force follows the activation dips --- the expected
smoothing of the series elastic element. The two outliers are
anatomically explained: rectus femoris was rerouted in the conversion
over the vastii's patella-tracking via point, and its moment-arm
polyfit deviates from OpenSim's path (the same coupler artifact noted
for the vastii); biceps femoris short head agrees in shape
($r^2 = 0.87$) but carries a scale offset. We conclude that the
converted plant reproduces OpenSim muscle force generation along
physiological trajectories to a degree sufficient for neural-control
studies, and that if tendon-elasticity phenomena (energy storage,
force smoothing) become material to a claim, the appropriate upgrade
is a compliant-tendon formulation rather than parameter retuning of
the rigid-tendon model.

\begin{figure}[h]
  \centering
  \includegraphics[width=0.9\textwidth]{muscle_force_compare.png}
  \caption{Normalized muscle force along the analyzed walk at matched
  OpenSim static-optimization activations (MuJoCo: rigid tendon;
  OpenSim: Thelen compliant tendon). Soleus illustrates the signature
  of the missing series compliance: OpenSim force is carried through
  mid-stance by tendon stretch while MuJoCo force follows the
  activation.}
  \label{fig:force-compare}
\end{figure}
